
\documentclass{article}
\usepackage{colortbl}
\usepackage{makecell}
\usepackage{multirow}
\usepackage{supertabular}

\begin{document}

\newcounter{utterance}

\centering \large Interaction Transcript for game `wordle\_withclue', experiment `high\_frequency\_words\_clue\_no\_critic', episode 5 with Qwen3.5{-}9B.
\vspace{24pt}

{ \footnotesize  \setcounter{utterance}{1}
\setlength{\tabcolsep}{0pt}
\begin{supertabular}{c@{$\;$}|p{.15\linewidth}@{}p{.15\linewidth}p{.15\linewidth}p{.15\linewidth}p{.15\linewidth}p{.15\linewidth}}
    \# & \multicolumn{2}{c}{Player} && \multicolumn{2}{c}{Game Master} \\
    \hline

    \theutterance \stepcounter{utterance}  
    & & & \multicolumn{4}{p{0.6\linewidth}}{
        \cellcolor[rgb]{0.9,0.9,0.9}{
            \makecell[{{p{\linewidth}}}]{
                \texttt{\tiny{[P1$\langle$GM]}}
                \texttt{You are a language wizard who likes to guess words by using the given rules.} \\
\\ 
\texttt{Welcome to Wordle! You have six attempts to guess the target word, a valid English word of five lowercase letters (a{-}z). Please use the tags "explanation:" and "guess:" to provide a concise explanation for each guess.} \\
\\ 
\texttt{To help you make an informed guess, you will receive a clue for the word, such as} \\
\texttt{clue: snowy white.} \\
\\ 
\texttt{Here is an example guess based on the clue:} \\
\texttt{explanation: In the fairy tail Snow White, the girl is killed because she eats a poisoned apple. And the word apple has 5 letters.} \\
\texttt{guess: apple} \\
\\ 
\texttt{After each guess, your answer will be validated, and you will receive feedback indicating which letters are correct (green), which letters are correct but in the wrong position (yellow), and which letters are incorrect (red). This feedback can be useful in determining which letters to include or exclude in your next guess.} \\
\\ 
\texttt{For example, the feedback for "apple" might be:} \\
\texttt{guess\_feedback: a<yellow> p<yellow> p<green> l<yellow> e<red>} \\
\\ 
\texttt{The explanation should contain details about how both the guess\_feedback and clue work together in tandem to formulate a new guess.} \\
\\ 
\texttt{Do not generate any other text.} \\
\\ 
\texttt{Let's begin with your first guess.} \\
\\ 
\texttt{clue: divine \_\_\_} \\
            }
        }
    }
     \\ \\

    \theutterance \stepcounter{utterance}  
    & \multicolumn{4}{p{0.6\linewidth}}{
        \cellcolor[rgb]{0.561,0.737,0.561}{
            \makecell[{{p{\linewidth}}}]{
                \texttt{\tiny{[P1$\rangle$GM]}}
                \texttt{explanation: The clue "divine \_\_\_" suggests a word related to divinity or a deity. A common five{-}letter word fitting this context is "godly," which means having the qualities of a god.} \\
\texttt{guess: godly} \\
            }
        }
    }
    & & \\ \\

    \theutterance \stepcounter{utterance}  
    & & & \multicolumn{2}{p{0.3\linewidth}}{
        \cellcolor[rgb]{0.9,0.9,0.9}{
            \makecell[{{p{\linewidth}}}]{
                \texttt{\tiny{[GM$|$GM]}}
                \texttt{attempts = 1} \\
\texttt{target\_word = right} \\
\texttt{target\_word\_clue = divine \_\_\_} \\
\texttt{guess = godly} \\
\texttt{guess\_feedback = g<yellow> o<red> d<red> l<red> y<red>} \\
            }
        }
    }
    & & \\ \\

    \theutterance \stepcounter{utterance}  
    & & & \multicolumn{4}{p{0.6\linewidth}}{
        \cellcolor[rgb]{0.9,0.9,0.9}{
            \makecell[{{p{\linewidth}}}]{
                \texttt{\tiny{[P1$\langle$GM]}}
                \texttt{guess\_feedback: g<yellow> o<red> d<red> l<red> y<red>} \\
\\ 
\texttt{Provide your response only in this format.} \\
\texttt{explanation: details} \\
\texttt{guess: word} \\
            }
        }
    }
     \\ \\

    \theutterance \stepcounter{utterance}  
    & \multicolumn{4}{p{0.6\linewidth}}{
        \cellcolor[rgb]{0.561,0.737,0.561}{
            \makecell[{{p{\linewidth}}}]{
                \texttt{\tiny{[P1$\rangle$GM]}}
                \texttt{explanation: The feedback for "godly" shows that 'g' is in the correct position (yellow), while 'o', 'd', 'l', and 'y' are incorrect (red). This eliminates 'o', 'd', 'l', and 'y' from the word. Since the clue is "divine \_\_\_", the word likely starts with 'g' and relates to divinity. With 'g' confirmed at the start and other letters ruled out, I need a word starting with 'g' that fits the divine context without using the eliminated letters. "grace" is a strong candidate as it relates to divine favor and fits the letter constraints.} \\
\texttt{guess: grace} \\
            }
        }
    }
    & & \\ \\

    \theutterance \stepcounter{utterance}  
    & & & \multicolumn{2}{p{0.3\linewidth}}{
        \cellcolor[rgb]{0.9,0.9,0.9}{
            \makecell[{{p{\linewidth}}}]{
                \texttt{\tiny{[GM$|$GM]}}
                \texttt{attempts = 2} \\
\texttt{target\_word = right} \\
\texttt{target\_word\_clue = divine \_\_\_} \\
\texttt{guess = grace} \\
\texttt{guess\_feedback = g<yellow> r<yellow> a<red> c<red> e<red>} \\
            }
        }
    }
    & & \\ \\

    \theutterance \stepcounter{utterance}  
    & & & \multicolumn{4}{p{0.6\linewidth}}{
        \cellcolor[rgb]{0.9,0.9,0.9}{
            \makecell[{{p{\linewidth}}}]{
                \texttt{\tiny{[P1$\langle$GM]}}
                \texttt{guess\_feedback: g<yellow> r<yellow> a<red> c<red> e<red>} \\
\\ 
\texttt{Provide your response only in this format.} \\
\texttt{explanation: details} \\
\texttt{guess: word} \\
            }
        }
    }
     \\ \\

    \theutterance \stepcounter{utterance}  
    & \multicolumn{4}{p{0.6\linewidth}}{
        \cellcolor[rgb]{0.561,0.737,0.561}{
            \makecell[{{p{\linewidth}}}]{
                \texttt{\tiny{[P1$\rangle$GM]}}
                \texttt{explanation: The feedback for "grace" indicates that 'g' is in the correct position (yellow) and 'r' is correct but in the wrong position (yellow), while 'a', 'c', and 'e' are incorrect (red). This confirms the word starts with 'g' and contains an 'r' somewhere else. Given the clue "divine \_\_\_", the word likely refers to a divine attribute or entity. With 'g' fixed at the start and 'r' misplaced, and 'a', 'c', 'e' eliminated, "grace" was a good guess but the 'r' is not in the second spot. Considering divine concepts, "grace" fits well but the position of 'r' is wrong. Let's try "grasp" or "great"? No, "great" has 'e' and 'a' which are red. How about "grind"? No. What about "glory"? 'l' and 'o' and 'y' were red in the first guess. Wait, 'g' is yellow (correct pos), 'r' is yellow (wrong pos). So the word is g \_ \_ \_ \_ with an r not in pos 2. Clue "divine \_\_\_". Could it be "grace" but r is in pos 3? No, feedback said r is yellow (wrong pos). So r is not in pos 2. Maybe "grace" is wrong. Let's think of words starting with G related to divine. "Godly" was tried. "Grace" was tried. What about "grace" feedback: g(y), r(y), a(r), c(r), e(r). So g is 1st. r is somewhere else. a,c,e are out. Clue: divine \_\_\_. Maybe "grace" is not the word. What about "grasp"? No. "grind"? No. "grove"? No. "grout"? No. "grime"? No. "grime" has i,m,e. e is red. "grime" has i,m. Maybe "grace" is close but r is in wrong spot. If r is yellow, it's in the word but not pos 2. So g \_ \_ r \_ or g \_ \_ \_ r. Divine \_\_\_? "grace" is divine favor. Maybe the word is "grace" but the feedback interpretation: g(yellow) means g is correct but wrong pos? Wait, standard Wordle: Green = correct pos, Yellow = correct letter wrong pos, Red = wrong letter. The prompt says: "g<yellow>". In the first guess "godly": g<yellow> o<red> d<red> l<red> y<red>. This means 'g' is in the word but NOT in position 1. But the clue is "div} \\
            }
        }
    }
    & & \\ \\

    \theutterance \stepcounter{utterance}  
    & & & \multicolumn{2}{p{0.3\linewidth}}{
        \cellcolor[rgb]{0.9,0.9,0.9}{
            \makecell[{{p{\linewidth}}}]{
                \texttt{\tiny{[GM$|$GM]}}
                \texttt{Error: The response should contain the 'guess:' keyword exactly once.} \\
            }
        }
    }
    & & \\ \\

    \theutterance \stepcounter{utterance}  
    & & & \multicolumn{2}{p{0.3\linewidth}}{
        \cellcolor[rgb]{0.9,0.9,0.9}{
            \makecell[{{p{\linewidth}}}]{
                \texttt{\tiny{[GM$|$GM]}}
                \texttt{game\_result = ABORT} \\
            }
        }
    }
    & & \\ \\

\end{supertabular}
}

\end{document}
